% Source: https://artofproblemsolving.com/community/c3h2669400_111th_post_solving_problems
% Stable identity: aops-topic-2669400-q20
By induction, $f^n(x)=a^nx+b\left(1+a+\cdots+a^{n-1}\right)$.
\begin{itemize}
\item $a=1$: $f^n(x)=x+nb$, which converges for all $x$ exactly when $b=0$.
\item $a\ne1$: $f^n(x)=a^n\left(x-c\right)+c$ with $c=\dfrac{b}{1-a}$ (the fixed point). If $|a|<1$ this tends to $c$ for every $x$. If $|a|\ge1$ (with $a\ne1$), then for $x\ne c$ the term $a^n(x-c)$ does not converge (it grows if $|a|>1$, and alternates between $\pm(x-c)$ if $a=-1$).
\end{itemize}
So the answer is $\boxed{|a|<1\text{ with any }b,\ \text{or }a=1,\ b=0}$.
